\section{関連研究}
\subsection{閉域的な物体検出モデル}
YOLO\cite{ref_YOLO}は画像を格子に分割し，各格子からボックスとクラスを1回の推論で予測する．
DETR\cite{ref_DETR}はTransformerによる集合予測として検出を定式化し，Non-Maximum Suppression（NMS）を用いずにボックスとクラスを出力する．
いずれも学習時に定めたクラス集合を前提としており，COCO\cite{ref_COCO}で学習したモデルは80クラス以外の語を出力できない．
また，学習データに存在しない物体は背景として扱われるため，検出そのものが行われないことがある．

\subsection{オープンボキャブラリ物体検出とVLM}
Contrastive Language-Image Pre-training（CLIP）\cite{ref_CLIP}は画像とテキストの対の対照学習により，任意の語に対するゼロショット分類を可能にした．
SigLIP\cite{ref_SigLIP}は，候補全体で正規化するソフトマックスの代わりに，画像とテキストの対ごとのシグモイドを用いる目的関数で学習され，語ごとに独立な確信度を与える．
OWL-ViT\cite{ref_OWLViT}やOWLv2\cite{ref_OWLv2}，Grounding DINO\cite{ref_GroundingDINO}は，カテゴリ名をテキストのクエリとして与えることで固定のクラス集合に依存しない検出を行う．
これらのモデルは検出とカテゴリの予測を単一のモデルで同時に行うため，候補語彙や粒度を変えるたびに検出を再実行する必要がある．
一方，OWLv2はテキストに依存しない物体らしさのスコアも出力し，未知の物体を含む領域の候補の生成に用いることができる．
本研究は，検出を1回に限り，候補語彙の変更をクロップの分類のみで行う点で，これらのモデルと異なる．
